% Solutions to Chapter 1's exercises, from lectures/L01/appendix/d_solutions.md. Every exercise in
% the chapter is worked on paper, so every one is answered.
%
% Exercise 10 a) answers the text IC74, as the book's exercise asks; see chapters/01/exercises.tex.

\answerchapter{c1:ch}

\answer{c1:ex:1}{Positionssystem}
\pt{a}Positionerna räknas från höger med början på noll:

\begin{narrowtable}{@{}cccc@{}}
  \thead{Siffra} & \thead{Position} & \thead{Vikt} & \thead{Bidrag} \\
  \midrule
  4 & 3 & $10^3 = 1000$ & 4000 \\
  0 & 2 & $10^2 = 100$ & 0 \\
  7 & 1 & $10^1 = 10$ & 70 \\
  2 & 0 & $10^0 = 1$ & 2 \\
\end{narrowtable}

\noindent Summan är $4000 + 0 + 70 + 2 = 4072$. Nollan bidrar inte med något, men den behövs ändå:
utan den skulle sjuan hamna på position 2 och vara värd 700.

\pt{b}Position 2 har vikten $5^2 = 25$. Då är
\[
  243_5 = 2 \cdot 25 + 4 \cdot 5 + 3 \cdot 1 = 50 + 20 + 3 = 73_{10}
\]

\pt{c}Vikten för position $n$ är $B^n$, och alla tal upphöjda till noll är 1. Ett mer vardagligt
sätt att se det: den högra positionen räknar ental, hur många \enquote{hela} det finns, och det är
samma sak i alla talsystem. Först när entalen tar slut, vid talbasen, börjar nästa position räkna.

\answer{c1:ex:2}{Från binärt till decimalt}
Summera vikterna för de positioner där det står en etta.

\pt{a}\code{1101} = 8 + 4 + 1 = \textbf{13}.
\pt{b}\code{1 0000} = \textbf{16}. En enda etta på position 4, alltså $2^4$.
\pt{c}\code{0110 0100} = 64 + 32 + 4 = \textbf{100}.
\pt{d}\code{1000 0001} = 128 + 1 = \textbf{129}.
\pt{e}\code{1111 1111} = 128 + 64 + 32 + 16 + 8 + 4 + 2 + 1 = \textbf{255}, det största 8-bitars
talet.
\pt{f}\code{1101 0110} = 128 + 64 + 16 + 4 + 2 = \textbf{214}.

\ptnote{Vanligt fel}Att börja viktraden från vänster med 1, alltså läsa talet baklänges.
\code{1101} blir då 1 + 2 + 8 = 11 i stället för 13. Viktraden börjar alltid med 1 längst till
\textbf{höger}.

\answer{c1:ex:3}{Från decimalt till binärt}
\pt{a}13 med vikttabellen: 8 ryms (kvar 5), 4 ryms (kvar 1), 2 ryms inte, 1 ryms (kvar 0). Svar:
\textbf{\code{1101}}. Kontroll: 8 + 4 + 1 = 13.

\pt{b}37: 32 ryms (kvar 5), 16 nej, 8 nej, 4 ja (kvar 1), 2 nej, 1 ja. Svar: \textbf{\code{10
0101}}, i en byte \code{0010 0101}. Kontroll: 32 + 4 + 1 = 37.

\pt{c}100: 64 ja (kvar 36), 32 ja (kvar 4), 16 nej, 8 nej, 4 ja (kvar 0), 2 nej, 1 nej. Svar:
\textbf{\code{110 0100}}, i en byte \code{0110 0100}. Jämför med övning~\ref{c1:ex:2} c).

\pt{d}200 med division med 2:

\begin{narrowtable}{@{}ccc@{}}
  \thead{Division} & \thead{Kvot} & \thead{Rest} \\
  \midrule
  200 / 2 & 100 & 0 \\
  100 / 2 & 50 & 0 \\
  50 / 2 & 25 & 0 \\
  25 / 2 & 12 & 1 \\
  12 / 2 & 6 & 0 \\
  6 / 2 & 3 & 0 \\
  3 / 2 & 1 & 1 \\
  1 / 2 & 0 & 1 \\
\end{narrowtable}

\noindent Nedifrån och upp: \textbf{\code{1100 1000}}. Kontroll: 128 + 64 + 8 = 200.

Lägg märke till att 200 är dubbelt så mycket som 100, och att \code{1100 1000} är \code{0110 0100}
flyttat ett steg åt vänster. Det är ingen slump, se övning~\ref{c1:ex:14}.

\pt{e}255: divisionerna ger rest 1 varje gång (255, 127, 63, 31, 15, 7, 3, 1), så svaret är
\textbf{\code{1111 1111}}. Det går också att se direkt: 255 = 256 -- 1, och $2^8 - 1$ är åtta
ettor.

\pt{f}1000:

\begin{narrowtable}{@{}ccc@{}}
  \thead{Division} & \thead{Kvot} & \thead{Rest} \\
  \midrule
  1000 / 2 & 500 & 0 \\
  500 / 2 & 250 & 0 \\
  250 / 2 & 125 & 0 \\
  125 / 2 & 62 & 1 \\
  62 / 2 & 31 & 0 \\
  31 / 2 & 15 & 1 \\
  15 / 2 & 7 & 1 \\
  7 / 2 & 3 & 1 \\
  3 / 2 & 1 & 1 \\
  1 / 2 & 0 & 1 \\
\end{narrowtable}

\noindent Nedifrån och upp: \textbf{\code{11 1110 1000}}, tio bitar. Kontroll: 512 + 256 + 128 + 64
+ 32 + 8 = 1000. Talet ryms alltså inte i en byte.

\answer{c1:ex:4}{Binärt och hexadecimalt}
Gruppera i fyror från höger, och byt varje grupp mot sin siffra.

\pt{a}\code{1010 1111} = \code{1010} \code{1111} = A F = \textbf{\code{0xAF}}.
\pt{b}\code{0011 1100} = 3 C = \textbf{\code{0x3C}}.
\pt{c}\code{1 1110} fylls ut till \code{0001 1110} = 1 E = \textbf{\code{0x1E}}. Den som grupperar
från vänster får \code{1111} \code{0} och svaret F0, som är ett helt annat tal: 240 i stället för
30.
\pt{d}\code{11 1110 1000} fylls ut till \code{0011 1110 1000} = 3 E 8 = \textbf{\code{0x3E8}}.
Jämför med övning~\ref{c1:ex:3} f): det är 1000.
\pt{e}\code{0x7F} = 7 F = \textbf{\code{0111 1111}}.
\pt{f}\code{0xC3} = C 3 = \textbf{\code{1100 0011}}.
\pt{g}\code{0x2A5} = 2 A 5 = \textbf{\code{0010 1010 0101}}, eller utan de inledande nollorna
\code{10 1010 0101}.

\answer{c1:ex:5}{Hexadecimalt och decimalt}
\pt{a}\code{0x2A} = 2 · 16 + 10 = \textbf{42}.
\pt{b}\code{0xFF} = 15 · 16 + 15 = 240 + 15 = \textbf{255}.
\pt{c}\code{0x100} = 1 · 256 + 0 + 0 = \textbf{256}. Ett mer än \code{0xFF}, och det första talet
som behöver tre hexadecimala siffror och nio bitar.
\pt{d}\code{0x3E8} = 3 · 256 + 14 · 16 + 8 = 768 + 224 + 8 = \textbf{1000}.
\pt{e}75 = 4 · 16 + 11, så siffrorna är 4 och B: \textbf{\code{0x4B}}. Via binärt: 75 = 64 + 8 + 2
+ 1 = \code{0100 1011} = 4B.
\pt{f}160 = 10 · 16 + 0: \textbf{\code{0xA0}}.
\pt{g}4095 = 4096 -- 1 = $2^{12} - 1$, alltså tolv ettor: \code{1111 1111 1111} =
\textbf{\code{0xFFF}}. Med division: 4095 / 16 = 255 rest 15, 255 / 16 = 15 rest 15, 15 / 16 = 0
rest 15. Tre femton, FFF.

\answer{c1:ex:6}{Bitar och talområden}
\pt{a}\leavevmode

\begin{narrowtable}{@{}ccc@{}}
  \thead{Bitar} & \thead{Antal värden} & \thead{Största värde} \\
  \midrule
  4 & $2^4 = 16$ & 15 \\
  8 & $2^8 = 256$ & 255 \\
  12 & $2^{12} = 4096$ & 4095 \\
  16 & $2^{16} = 65\,536$ & 65\,535 \\
\end{narrowtable}

\pt{b}$2^9 = 512$ räcker inte, men $2^{10} = 1024$ gör det. Alltså \textbf{10 bitar}.

\pt{c}Nivåerna 0 till 200 är \textbf{201} olika värden, inte 200: noll räknas också. Sju bitar ger
128 värden, vilket inte räcker; åtta bitar ger 256. Alltså \textbf{8 bitar}, och 256 -- 201 =
\textbf{55} kombinationer blir oanvända.

Det är lätt att svara 200 värden, och i det här fallet ändrar det inte antalet bitar. Men i en tank
som ska mäta 0 till 128 gör det skillnad: det är 129 värden, som kräver åtta bitar, medan den som
räknar till 128 värden drar slutsatsen att sju bitar räcker.

\pt{d}LSB är \code{1}, så talet är \textbf{udda}: alla andra vikter är jämna. MSB är \code{1}, så
talet är \textbf{minst 128}: MSB ensam är värd 128. (Talet är 183.)

\answer{c1:ex:7}{Binär addition}
\pt{a}58 + 22:

\begin{textdrawing}
  minne:   0111 1100
           0011 1010     (58)
         + 0001 0110     (22)
         -----------
           0101 0000     (80)
\end{textdrawing}

\noindent\code{0101 0000} = 64 + 16 = 80. Stämmer.

\pt{b}127 + 1:

\begin{textdrawing}
  minne:   1111 1110
           0111 1111     (127)
         + 0000 0001     (1)
         -----------
           1000 0000     (128)
\end{textdrawing}

\noindent Minnessiffran vandrar genom alla sju ettorna, precis som 999 + 1 i decimalt.

\pt{c}200 + 100:

\begin{textdrawing}
  minne: 1 1000 0000
           1100 1000     (200)
         + 0110 0100     (100)
         -----------
         1 0010 1100
\end{textdrawing}

\pt{d}Summan i c) är 300 och ryms inte i åtta bitar. I registret blir \code{0010 1100} kvar, alltså
\textbf{44}, och den nionde biten, minnessiffran ut ur bit~7, hamnar utanför. 300 -- 256 = 44: det
som försvann var exakt 256, värdet av den nionde biten. I en mikrodator sparas den biten i flaggan
\emph{carry}, så att programmet kan upptäcka att summan inte fick plats.

\answer{c1:ex:8}{BCD}
\pt{a}47 i BCD: 4 = \code{0100} och 7 = \code{0111}, alltså \textbf{\code{0100 0111}}. Som binärt
tal är 47 = 32 + 8 + 4 + 2 + 1 = \textbf{\code{0010 1111}}. Två olika bitmönster för samma tal.

\pt{b}2026 i BCD: \textbf{\code{0010 0000 0010 0110}}, fyra bitar per siffra, 16 bitar. Som binärt
tal behövs bara 11 bitar (2026 = \code{111 1110 1010}), så BCD kostar fler bitar.

\pt{c}\code{1001 0011} i BCD är 9 och 3, alltså \textbf{93}. Som binärt tal är samma mönster 128 +
16 + 2 + 1 = \textbf{147}.

\pt{d}\textbf{Nej.} Den högra nibblen \code{1010} är 10, och i BCD finns bara siffrorna 0--9. En
BCD-krets som får \code{1010} har fått något som inte är en siffra, och vad den gör med det är
odefinierat.

\answer{c1:ex:9}{Graykod}
\pt{a}Tvåbitarskoden och dess spegelbild, med 0 respektive 1 framför:

\begin{console}
0 00   0 01   0 11   0 10   1 10   1 11   1 01   1 00
\end{console}

\noindent Alltså \code{000, 001, 011, 010, 110, 111, 101, 100}, samma som tabellen i
\secref{c1:b3}.

\pt{b}Binärt, från \code{111} till \code{000}: \textbf{alla tre} bitarna ändras. Graykod, från
\code{100} till \code{000}: \textbf{en} bit. Graykoden håller egenskapen hela vägen runt, också när
räknaren slår om från det sista talet till det första.

\pt{c}Binärt är läge 1 \code{001} och läge 2 \code{010}: både bit~1 och bit~0 ändras. Hinner bit~0
ändras först läser styrsystemet ett ögonblick \code{000}, alltså läge 0; hinner bit~1 ändras först
läser det \code{011}, alltså läge 3. Givaren kan alltså tillfälligt rapportera ett läge som axeln
aldrig var i.

Med Graykod är läge 1 \code{001} och läge 2 \code{011}: bara bit~1 ändras. Under övergången läses
antingen \code{001} eller \code{011}, alltså läge 1 eller läge 2, och båda är rimliga svar på var
axeln är just då.

\answer{c1:ex:10}{ASCII}
\pt{a}\code{I} = \code{0x49} (\code{A} är \code{0x41}, så \code{I}, åtta bokstäver senare, är
\code{0x49}), \code{C} = \code{0x43}, \code{7} = \code{0x37}, \code{4} = \code{0x34}. Texten är
\textbf{\code{0x49 0x43 0x37 0x34}}.

Siffertecknen 7 och 4 blir \code{0x37} och \code{0x34}, inte \code{0x07} och \code{0x04}. Det är det
vanligaste felet i den här uppgiften.

\pt{b}\code{A} är \code{0x41}, så \code{0x48} är sju bokstäver senare: B, C, D, E, F, G,
\textbf{H}. \code{0x65} är \textbf{e}, och \code{0x6A} är fem bokstäver efter e, \textbf{j}. Texten
är \textbf{\code{Hej}}.

\pt{c}\code{a} är \code{0x61} och \code{A} är \code{0x41}, en skillnad på \textbf{\code{0x20}} = 32.
Binärt är \code{A} = \code{0100 0001} och \code{a} = \code{0110 0001}: de skiljer sig bara i
\textbf{bit~5}, som har vikten 32.

\pt{d}Siffertecknet \code{7} är \code{0x37}. Programmet drar bort \code{0x30} och får \textbf{7}. Åt
andra hållet lägger det till \code{0x30}: 3 + \code{0x30} = \code{0x33}, som är tecknet
\textbf{\code{3}}. Det fungerar bara för en siffra 0--9; tal med flera siffror måste delas upp
siffra för siffra först.

\answer{c1:ex:11}{7-segmentkod}
\pt{a}\code{E} ritas med a, d, e, f och g:

\begin{narrowtable}{@{}ccccccc@{}}
  g & f & e & d & c & b & a \\
  \midrule
  1 & 1 & 1 & 1 & 0 & 0 & 1 \\
\end{narrowtable}

\noindent Binärt \code{111 1001}, i en byte \code{0111 1001}, hexadecimalt \textbf{\code{0x79}}.

\pt{b}\code{0x6D} = \code{0110 1101}. Med g f e d c b a för bit 6 till 0: g = 1, f = 1, e = 0,
d = 1, c = 1, b = 0, a = 1. Segmenten a, c, d, f och g lyser, vilket är siffran \textbf{5}. Jämför
med figur~\ref{c1:fig:seven_segment}.

\pt{c}Siffror som inte använder g påverkas inte: \textbf{0, 1 och 7}. Den allvarliga förväxlingen
är \textbf{8}: utan g lyser a, b, c, d, e och f, vilket är exakt samma segment som \textbf{0}. En 8
och en 0 går alltså inte att skilja åt.

Övriga siffror som använder g, 2, 3, 4, 5, 6 och 9, blir ofullständiga figurer som inte liknar
någon siffra. De är fel, men felet syns. Det är 8:an som är farlig, eftersom den ser rimlig ut: en
display som visar 0 grader när det är 8 grader väcker ingen misstanke.

\answer{c1:ex:12}{Paritetsbit}
\pt{a}\leavevmode

\begin{narrowtable}{@{}ccccc@{}}
  \thead{Tecken} & \thead{ASCII (7 bitar)} & \thead{Antal ettor} & \thead{Paritetsbit} &
    \thead{Byte som skickas} \\
  \midrule
  \code{B} & \code{100 0010} & 2 & 0 & \code{0100 0010} \\
  \code{C} & \code{100 0011} & 3 & 1 & \code{1100 0011} \\
  \code{G} & \code{100 0111} & 4 & 0 & \code{0100 0111} \\
\end{narrowtable}

\pt{b}\code{1100 0001} har tre ettor, ett udda antal, så \textbf{ett fel har inträffat}. Men
mottagaren kan inte veta vilken bit som vändes. Det kan ha varit paritetsbiten (då var tecknet
\code{A}, \code{0100 0001}), eller någon av de andra bitarna. Allt mottagaren kan göra är att be om
tecknet igen.

\pt{c}Ett exempel: \code{A} skickas som \code{0100 0001}. Två störningar vänder bit~0 och bit~1, och
mottagaren får \code{0100 0010}. Det har två ettor, ett jämnt antal, så paritetskontrollen godkänner
byten. Mottagaren tar emot \code{B} i stället för \code{A} utan att märka något. Varje dubbelfel
ändrar antalet ettor med två, noll eller minus två, och antalet förblir jämnt.

\answer{c1:ex:13}{Kontroll: för hand, och sedan med kalkylatorn}
\pt{a}\leavevmode
\begin{itemize}
  \item 173 = 128 + 32 + 8 + 4 + 1 = \textbf{\code{1010 1101}} = \textbf{\code{0xAD}}.
  \item \code{0x3E8} = 3 · 256 + 14 · 16 + 8 = \textbf{1000}, binärt \textbf{\code{11 1110 1000}}.
  \item \code{0b1100 1010} = 128 + 64 + 8 + 2 = \textbf{202} = \textbf{\code{0xCA}}.
\end{itemize}

\pt{b}Kalkylatorn ska visa exakt dessa värden. Om den inte gör det är det nästan alltid ett av tre
fel: en vikt som hoppats över i summeringen, resterna i divisionen lästa i fel ordning, eller
bitarna grupperade från vänster i stället för från höger. Hitta vilket, så att du inte gör det
igen.

\pt{c}I läget \code{QWORD} (64 bitar) visar kalkylatorn \code{-1} som \code{FFFF FFFF FFFF FFFF},
och binärt som 64 ettor. I läget \code{BYTE} visar den \code{FF} och \code{1111 1111}.

\pt{d}\code{1111 1111} är som positivt tal \textbf{255}. Ett skäl till att datorn skriver $-1$ så:
lägg till 1, och vad händer? Enligt \secref{c1:a7} blir 255 + 1 = 256, som i åtta bitar är
\code{0000 0000}, alltså 0. Ett tal som blir noll när man lägger till ett beter sig precis som $-1$.
Det är idén bakom \textbf{tvåkomplement}, som är ämnet för kapitel~\ref{c8:ch}. Om du gissade något
i den riktningen har du redan förstått det viktigaste.

\answer{c1:ex:14}{Att lägga till en nolla}
\pt{a}\code{0b10110} = 16 + 4 + 2 = \textbf{22}, och \code{0b101100} = 32 + 8 + 4 = \textbf{44}.
Varje nolla som läggs till fördubblar talet.

\pt{b}När alla bitar flyttas ett steg åt vänster hamnar varje etta på en position som är värd
dubbelt så mycket. \textbf{Värdet fördubblas}, precis som ett decimalt tal blir tio gånger större
när man lägger till en nolla. (I ett register med åtta bitar gäller det bara så länge ingen etta
trillar ut till vänster, se övning~\ref{c1:ex:7} d).)

\pt{c}Värdet \textbf{halveras, och en eventuell rest försvinner}: \code{0b1011} (11) blir
\code{0b101} (5), och \code{0b1010} (10) blir också \code{0b101} (5). Biten som tas bort är LSB, som
avgör om talet är udda, så den försvunna halvan är just resten vid division med 2.

\pt{d}Talet blir \textbf{16 gånger större}, eftersom talbasen är 16: \code{0x2A} (42) blir
\code{0x2A0} = 2 · 256 + 10 · 16 = 672 = 42 · 16. En hexadecimal siffra är fyra bitar, så att lägga
till en hexadecimal nolla är detsamma som att flytta alla bitar fyra steg åt vänster: $2^4 = 16$.
